\subsection{FRM Exam 2001: Rank portfolios from least risky to most risky}

\textbf{Enunciado}

\emph{Rank the following portfolios from least risky to most risky. Assume 252 trading days a year and there are 5 trading days per week.}
\[
\begin{array}{c c c c}
\toprule
\text{Portfolio} & \text{VaR} & \text{Holding period (days)} & \text{Confidence}\\
\midrule
1 & 10 & 5  & 99\\
2 & 10 & 5  & 95\\
3 & 10 & 10 & 99\\
4 & 10 & 10 & 95\\
5 & 10 & 15 & 99\\
6 & 10 & 15 & 95\\
\bottomrule
\end{array}
\]
\begin{enumerate}[label=\alph*.]
\item \hl{5, 3, 6, 1, 4, 2}
\item 3, 4, 1, 2, 5, 6
\item 5, 6, 1, 2, 3, 4
\item 2, 1, 5, 6, 4, 3
\end{enumerate}

\subsubsection*{Solución}

Bajo el supuesto delta-normal, para un mismo portafolio:
\[
\VaR_\alpha(h)\approx z_\alpha\,\sigma\,\sqrt{h},
\]
donde \(\sigma\) es la volatilidad diaria del P\&L, es un parámetro comparable entre portafolios.
Luego:
\[
\sigma \approx \frac{\VaR}{z_\alpha\sqrt{h}}.
\]
Usamos cuantiles de una cola: \(z_{0.99}\approx 2.33\), \(z_{0.95}\approx 1.645\).

\[
\sigma_1\approx \frac{10}{2.33\sqrt5}\approx 1.92,\qquad
\sigma_2\approx \frac{10}{1.645\sqrt5}\approx 2.72,
\]
\[
\sigma_3\approx \frac{10}{2.33\sqrt{10}}\approx 1.36,\qquad
\sigma_4\approx \frac{10}{1.645\sqrt{10}}\approx 1.92,
\]
\[
\sigma_5\approx \frac{10}{2.33\sqrt{15}}\approx 1.11,\qquad
\sigma_6\approx \frac{10}{1.645\sqrt{15}}\approx 1.57.
\]

Menor \(\sigma\) implica menor riesgo subyacente:
\[
\sigma_5<\sigma_3<\sigma_6<\sigma_1\approx \sigma_4<\sigma_2
\;\Rightarrow\;
5,3,6,1,4,2.
\]











